%% kapa: exam=1 round=23 year=2012 period=1 form=A subject=civil_law subject_ko=민법 seq=5 number=5
%% source: Q-Net 감정평가사 기출문제 (cst00309) artlSeq=5052578 file=raw/1008/5052578/2012년 제23회 감정평가사 제1차시험 기출문제(1교시 A형).pdf
%% conversion: kordoc 4.24.51 -> tools/gen1.py (auto); review=proofread-sonnet; flags=none
\begin{question}{5}
소멸시효 기간의 기산점에 관한 내용으로 옳지 않은 것은? (다툼이 있으면 판례에 의함)
\begin{choices}
  \item 동시이행의 항변권이 붙어 있는 채권은 이행기가 도래한 때부터 소멸시효가 진행한다.
  \item 부작위를 목적으로 하는 채권은 위반행위를 한 때로부터 소멸시효가 진행한다.
  \item 불확정기한부 채권은 채무자가 그 기한의 도래를 안 때부터 소멸시효가 진행한다.
  \item 기한을 정하고 있지 않은 채권은 그 채권이 발생한 때부터 소멸시효가 진행한다.
  \item 이행불능으로 인한 전보배상청구권의 소멸시효는 이행불능시부터 진행된다.
\end{choices}
\end{question}
\answer{3}
